\subsection*{Question 4.1 - Experimental amplitude Bode diagram of the closed-loop pitch}

The goal is to measure the amplitude of the closed-loop pitch response $\theta(t)$ to a sinusoidal pitch reference, for several frequencies in $\omega \in [1, 20]$ rad/s.

\paragraph{Step 1 - Model modification.}
Starting from \texttt{ARDroneHoverPitch2023.slx} (saved as a separate model, as recommended in the handout), the pitch reference is generated by a \emph{Sine Wave} block whose parameters are the workspace variables \texttt{pitchAmp} (amplitude) and \texttt{pitchFreq} (frequency, in rad/s).
% TODO: confirm that the "Pitch Source" manual switch is the one that selects the Sine Wave block in the saved model.
A manual switch selects between the sinusoid and the baseline (hover) pitch reference, so the drone takes off and stabilizes with the baseline controller, and the test signal is only applied once the drone is hovering.
Two \emph{To Workspace} blocks, with the same sample time $T_s = 0.03$ s, store the input and the output of the experiment:
\begin{itemize}
\item \texttt{thetaref\_data}: the pitch reference (input), which is the raw command in $[-1, 1]$ sent to the drone;
\item \texttt{theta\_data}: the measured pitch angle $\theta$ in rad (output).
\end{itemize}
Both are saved with the Structure With Time format and with the same sample time and decimation, so that they have the same length.

\paragraph{Step 2 - Experiment procedure.}
For each frequency, the following steps were repeated:
\begin{enumerate}
\item set \texttt{pitchFreq} and \texttt{pitchAmp} in the workspace;
\item take off and let the drone hover with the test signal switched off;
\item switch the test signal on and wait for the transient to die out, so that the response is a steady-state sinusoid;
\item switch the test signal off, land, stop the simulation, and save the run (\texttt{thetaref\_data} and \texttt{theta\_data}) into a file named after $\omega$ and the amplitude.
\end{enumerate}
The amplitude is increased with the frequency.
At low frequency a small amplitude is enough to get a clearly visible response, and it keeps the pitch angle small, where the linear model is valid.
At high frequency the closed-loop gain is much smaller, so a larger reference amplitude is needed to keep the output above the sensor noise.
All the commands stay well below the saturation limit of $1$.

\paragraph{Step 3 - From the recorded signal to the amplitude.}
Each recorded run contains the take-off, the hover, the test sinusoid and the landing, so the amplitude cannot be read from the whole run.
The processing has three stages.

\emph{1. Selection of the steady-state window.}
The height (hover at about $0.75$ m) and the amplitude of $\theta$ over time show when the test signal is switched on: the response grows from zero to a stable plateau.
For each run, a window after that transient and before landing was chosen, and it was cut to an integer number of periods of the sinusoid.
Only this window is used.

\emph{2. Removal of everything that is not the test sinusoid.}
In the window, the measured pitch is
\begin{equation*}
\theta(t) = A_{res}\sin(\omega t + \varphi) + c + n(t) ,
\end{equation*}
where $\omega$ is known (it is the frequency commanded), $c$ is a constant offset (the hover trim, so the mean of $\theta$ is not zero), and $n(t)$ is noise and other disturbances.
The mean of the signal gives only $c$, because the mean of a sinusoid over an integer number of periods is zero, so it cannot give the amplitude.
Reading the peaks of the plot gives a value that is biased upwards, since the noise adds to the peaks, and it uses only a few samples.
Instead, the unknowns are found by least squares from all the samples of the window, using the equivalent form
\begin{equation*}
\theta(t) = a\sin(\omega t) + b\cos(\omega t) + c ,
\end{equation*}
which is linear in $(a, b, c)$.
Over an integer number of periods, the sine, the cosine and the constant are orthogonal, so $a$ and $b$ are just the correlations of $\theta$ with $\sin(\omega t)$ and $\cos(\omega t)$ (like a single-frequency Fourier coefficient).
Noise at other frequencies averages out, and the offset is absorbed by $c$.

\emph{3. Amplitude and gain.}
The amplitude of the response is $A_{res} = \sqrt{a^2 + b^2}$, and the amplitude ratio is
\begin{equation*}
|G(j\omega)| = \frac{A_{res}}{A_{ref}}, \qquad |G(j\omega)|_{dB} = 20\log_{10}|G(j\omega)| ,
\end{equation*}
where $A_{ref}$ is the amplitude of the commanded sinusoid.
An equivalent (and simpler) check is $A_{res} \approx \sqrt{2}\,\theta_{rms}$ with the offset removed, which gives the same result when the noise is small.

The estimates are stable within each window: the amplitudes of sub-blocks of the same window differ by at most about $5\%$, and the standard error of the least-squares amplitude is at most about $0.0015$ rad, that is, about $5\%$ at $20$ rad/s and below $2\%$ at the other frequencies.

\paragraph{Step 4 - Measurements.}
Table~\ref{tab:pitch-bode} gives, for each frequency, the amplitude of the reference, the amplitude of the steady-state response, and the gain.
\begin{table}[ht]
\centering
\begin{tabular}{rrrcrr}
\hline
$\omega$ [rad/s] & $A_{ref}$ & $A_{res}$ [rad] & window [s] & $|G|$ & $|G|$ [dB] \\
\hline
1  & 0.1 & 0.0418 & 14 - 45.4   & 0.418 & $-7.6$  \\
5  & 0.2 & 0.0994 & 17 - 29.6   & 0.497 & $-6.1$  \\
10 & 0.3 & 0.0738 & 12 - 18.9   & 0.246 & $-12.2$ \\
15 & 0.4 & 0.0502 & 12.5 - 20.9 & 0.126 & $-18.0$ \\
20 & 0.5 & 0.0304 & 9.5 - 14.8  & 0.061 & $-24.3$ \\
\hline
\end{tabular}
\caption{Amplitude response of the closed-loop pitch obtained from the real runs.}
\label{tab:pitch-bode}
\end{table}

\paragraph{Step 5 - Amplitude diagram and comparison with the fitted model.}
Table~\ref{tab:pitch-bode-fit} compares the measured gain in dB with the second-order model with a zero identified in Question 4.2.
\begin{table}[ht]
\centering
\begin{tabular}{rrrr}
\hline
$\omega$ [rad/s] & Measured [dB] & Model [dB] & Error [dB] \\
\hline
1  & $-7.6$  & $-7.6$  & $0.0$  \\
5  & $-6.1$  & $-6.1$  & $0.0$  \\
10 & $-12.2$ & $-12.1$ & $-0.1$ \\
15 & $-18.0$ & $-18.8$ & $+0.8$ \\
20 & $-24.3$ & $-23.1$ & $-1.2$ \\
\hline
\end{tabular}
\caption{Measured and fitted gain of the closed-loop pitch. Error = measured $-$ model.}
\label{tab:pitch-bode-fit}
\end{table}
The gain is approximately constant (about $-7.6$ to $-6.1$ dB) up to $5$ rad/s, with a slight increase around $5$ rad/s, and then it decreases with the frequency, by about $18$ dB between $5$ and $20$ rad/s.
This is the behaviour of a low-pass system with a mild resonance.
The ratio is in rad per unit of command, so the low-frequency gain is not $0$ dB.
The model follows the data within $0.1$ dB up to $10$ rad/s, and differs by $0.8$ to $1.2$ dB at $15$ and $20$ rad/s.
% TODO: insert the Bode plot (fitted curve and experimental points) generated by fit_pitch_bode.m here.
